\documentclass[11pt]{amsart}

\usepackage[T1]{fontenc}
\usepackage{lmodern}
\usepackage{microtype}
\usepackage{amsmath,amssymb}
\usepackage[hidelinks]{hyperref}
\hypersetup{
  pdftitle={The Wagner graph refutes an El Zein--Mortada local-girth conjecture},
  pdfsubject={$S$-packing colorings of subcubic graphs},
  pdfkeywords={$S$-packing coloring, Wagner graph, local girth}
}

\newtheorem{theorem}{Theorem}
\newcommand{\Z}{\mathbb Z}
\newcommand{\dist}{\operatorname{dist}}

\title[A Wagner-graph counterexample]
      {The Wagner graph refutes an El Zein--Mortada local-girth conjecture}

\date{}

\subjclass[2020]{05C15}
\keywords{$S$-packing coloring, subcubic graph, Wagner graph, local girth}

\begin{document}

\begin{abstract}
El Zein and Mortada list as conjectured that every $(3,3)$-saturated
subcubic graph $G$ with $g_3(G)=4$ is $(1,2,2,2,2)$-packing colorable.
The Wagner graph is a counterexample. Every vertex has local girth $4$,
and the graph has diameter $2$ and independence number $3$; hence a
$(1,2,2,2,2)$-packing coloring can cover at most $3+4=7$ of its eight
vertices. We also prove that every subcubic graph on at most seven
vertices is $(1,2,2,2,2)$-packing colorable, so order eight is minimum.
Earlier computation of Gastineau and Togni already implied the existence
of such an obstruction. The present contribution is a standard explicit
witness and an elementary proof of minimum order.
\end{abstract}

\maketitle

Throughout, graphs are finite and simple, and vertices in different
components are taken to be at infinite distance. For a nondecreasing
sequence $S=(s_1,\ldots,s_k)$, an \emph{$S$-packing coloring} of a graph
$G$ is a partition
\[
  V(G)=V_1\mathbin{\dot\cup}\cdots\mathbin{\dot\cup}V_k
\]
such that $\dist_G(u,v)>s_i$ for all distinct $u,v\in V_i$.
The local girth $g(v)$ is the length of a shortest cycle containing $v$,
with $g(v)=\infty$ if no such cycle exists, and
\[
  g_3(G)=\max\{g(v):d_G(v)=3\}.
\]
Following El Zein and Mortada, a degree-three vertex is \emph{heavy} if
all its neighbors have degree three, and a subcubic graph is
$(3,k)$-saturated if every heavy vertex has at most $k$ heavy neighbors.
Thus every subcubic graph is $(3,3)$-saturated.

In version~1 of \cite{ElZeinMortadaImpact}, Table~1 lists the following
assertion in its conjectured column; it is the $g_3(G)=4$ case of the
first question in \cite[Problem~1]{ElZeinMortadaStep}:
\begin{quote}
\emph{Every $(3,3)$-saturated subcubic graph $G$ with $g_3(G)=4$ is
$(1,2,2,2,2)$-packing colorable.}
\end{quote}
We refer to this assertion as the El Zein--Mortada local-girth conjecture.

Let $W$ be the Wagner graph, with vertex set $\Z_8$ and edge set
\[
 E(W)=\bigl\{\{i,i+1\}:i\in\Z_8\bigr\}
      \cup
      \bigl\{\{i,i+4\}:i=0,1,2,3\bigr\},
\]
where all additions are modulo $8$.

\begin{theorem}\label{thm:main}
The graph $W$ is a $(3,3)$-saturated cubic graph with $g_3(W)=4$ that is
not $(1,2,2,2,2)$-packing colorable. Moreover, every subcubic graph of
order at most seven is $(1,2,2,2,2)$-packing colorable. Consequently,
$W$ is a minimum-order counterexample to the El Zein--Mortada local-girth
conjecture and to the first question in
\cite[Problem~1]{ElZeinMortadaStep}.
\end{theorem}

\begin{proof}
The graph $W$ is cubic, hence $(3,3)$-saturated. For every $i\in\Z_8$,
\[
  i,\ i+1,\ i+5,\ i+4
\]
induces a $4$-cycle. The neighbors $i-1,i+1,i+4$ of $i$ are pairwise
nonadjacent, so no triangle contains $i$. Thus every vertex has local
girth exactly $4$, and $g_3(W)=4$.

Every two vertices are at distance at most two. Cyclic differences
$1$ and $4$ give edges, difference $2$ is realized along the $8$-cycle,
and for difference $3$ one may use the path
$i,i+4,i+3$. Since $0$ and $2$ are nonadjacent,
$\operatorname{diam}(W)=2$.

The set $\{0,2,5\}$ is independent. An independent set of order four
in the spanning cycle $C_8$ must be one of
\[
  \{0,2,4,6\},\qquad \{1,3,5,7\},
\]
and each contains antipodal edges of $W$. Hence $\alpha(W)=3$.
In a $(1,2,2,2,2)$-packing coloring, the color-$1$ class therefore has
at most three vertices. Each of the four color-$2$ classes has at most
one vertex because $\operatorname{diam}(W)=2$. Such a coloring covers
at most $3+4=7<8$ vertices, a contradiction.

It remains to prove minimum order. Colors may be reused in different
components, so it suffices to color each connected component. A
component of order at most four can be colored by assigning its vertices
distinct color-$2$ classes. Let $G$ be connected with
$5\le |V(G)|\le7$. If $\Delta(G)\le2$, then $\chi(G)\le3$. If
$\Delta(G)=3$, Brooks' theorem also gives $\chi(G)\le3$, since $G$ is
neither $K_4$ nor an odd cycle. A largest class in a proper
$3$-coloring has at least $\lceil |V(G)|/3\rceil$ vertices. Assign that
independent class color $1$ and assign the remaining vertices distinct
color-$2$ classes. This is possible because
\[
  |V(G)|-\left\lceil\frac{|V(G)|}{3}\right\rceil\le4
  \qquad (5\le |V(G)|\le7).
\]
Thus every subcubic graph of order at most seven is
$(1,2,2,2,2)$-packing colorable.
\end{proof}

\section*{Priority and scope}

Gastineau and Togni's exhaustive table records two cubic graphs of order
$8$ that are $(1,2,2,2,2,2)$-packing colorable but not
$(1,2,2,2,2)$-packing colorable \cite[Table~1]{GastineauTogni}.
The Wagner graph is one of them: the classes
\[
  \{0,2,5\},\ \{1\},\ \{3\},\ \{4\},\ \{6\},\ \{7\}
\]
give a $(1,2,2,2,2,2)$-packing coloring, while
Theorem~\ref{thm:main} excludes $(1,2,2,2,2)$.

Their computation already implied a negative answer to the later
local-girth question. Indeed, if a vertex of an order-eight cubic graph
belonged to no triangle or $4$-cycle, its radius-two ball would contain
\[
  1+3+3\cdot2=10
\]
distinct vertices. The absence of triangles makes the first and second
layers disjoint, while the absence of $4$-cycles makes the six vertices
in the second layer distinct across the three branches. Thus every
vertex of an order-eight cubic graph has local girth at most $4$.

Only the $(1,2,2,2,2)$ assertion is refuted. The Wagner graph is
$(1,1,2,3)$-packing colorable, for example with color classes
\[
  \{0,2,5\},\qquad \{1,3,6\},\qquad \{4\},\qquad \{7\}.
\]
It therefore does not address the second question in
\cite[Problem~1]{ElZeinMortadaStep} or Conjecture~1 of that paper.

\begin{thebibliography}{9}

\bibitem{ElZeinMortadaStep}
A. El Zein and M. Mortada,
\emph{A step toward an $S$-packing coloring conjecture for subcubic graphs},
Australas. J. Combin. \textbf{95}(3) (2026), 459--467,
\url{https://ajc.maths.uq.edu.au/pdf/95/ajc_v95_p459.pdf}.

\bibitem{ElZeinMortadaImpact}
A. El Zein and M. Mortada,
\emph{Impact of local girth on the $S$-packing coloring of $k$-saturated
subcubic graphs},
\href{https://arxiv.org/abs/2603.25113v1}{arXiv:2603.25113v1} (2026).

\bibitem{GastineauTogni}
N. Gastineau and O. Togni,
\emph{$S$-packing colorings of cubic graphs},
Discrete Math. \textbf{339}(10) (2016), 2461--2470,
\href{https://doi.org/10.1016/j.disc.2016.04.017}
{doi:10.1016/j.disc.2016.04.017}.

\end{thebibliography}

\end{document}